%%%%%%%%%%%%%%%%%%%%%%%%%%%%%%%%%%%%%%%%%%%%%%%%%%%%%%%%%%%%%
\section{模型假设}

为简化问题并聚焦核心物理，本文做出以下假设。每条假设均给出合理性依据，并注明偏离真实情况时的处理方向。

\begin{enumerate}[itemindent=2em]
\item \textbf{双光束近似（A1）：} 只计表面反射与衬底界面一次反射两束光的干涉，忽略高阶往返反射束。\\
\textbf{依据：} 衬底界面反射系数幅值 $\rho_{12}\approx0.165$，高阶束以 $q=\rho_{12}|r_{01}|\approx0.072$ 逐级衰减，二阶项贡献约 7.2\%。\\
\textbf{去向：} 问题3 将判定并量化多光束效应的显著性，对 Si 样品启用 Airy 多光束模型。

\item \textbf{外延层折射率为实常数（A2）：} $n_1=2.55$，不随波数变化。\\
\textbf{依据：} 文献红外测量值（Goldberg 2001, 4H-SiC）；避开剩余射线带（$\sigma<1100$\,cm$^{-1}$）后中红外波段 $n$ 近乎平坦。数据实验（EXP-2）证实参数化色散与 $d$ 简并不可辨识。\\
\textbf{去向：} 色散导致的相位啁啾已计入系统误差区间 $[8.0,8.9]\,\mu$m；问题3 的 v2.0 模型通过衬底 Drude 频变相位部分补偿。

\item \textbf{衬底界面为常数反射率与常数有效相位（A3）：} $\rho_{12}$ 与 $\varphi_{12}$ 在测量波段内为常数。\\
\textbf{依据：} 数据实验（EXP-4）尝试 Drude 自由拟合，结果 $\gamma$ 误差棒 $\pm3700\%$、$d$ 漂移 $+31\%$，按预注册判据取常数模型并声明。\\
\textbf{去向：} v2.0 改用固定文献参数的 Drude 模型，释放了常数相位假设。

\item \textbf{衬底为光学厚半空间（A4）：} 无底面回波，衬底可视为半无限介质。\\
\textbf{依据：} 重掺 SiC 衬底自由载流子吸收强，红外穿透深度远小于衬底厚度（SEMI MF95 与 Thermo Fisher AN50639 的设定）。

\item \textbf{界面锐利无过渡层（A5）：} 外延层/衬底界面为理想阶跃折射率界面。\\
\textbf{依据：} 题目设定的单界面结构。\\
\textbf{去向：} 工艺中实际存在掺杂过渡层，问题3 的讨论中提及。

\item \textbf{平面波入射、样品横向均匀（A6）：} 红外光斑尺寸相对条纹间距尺度较小，忽略横向不均匀性。\\
\textbf{局限：} 晶圆面内均匀性未验证（无重复测量数据），不排除局部厚度梯度。

\item \textbf{偏振不敏感（A7）：} 探测器对 s/p 偏振等权响应，可用标量 Fresnel 系数。\\
\textbf{依据：} 数据实验（EXP-3）表明 s/p 四种设定 $\chi^2$ 差 $<4\%$，统计不可判别；近法线入射（$10\degree$--$15\degree$）下 s/p 差异仅 $1\sim3\%$。

\item \textbf{噪声平稳且点间独立（A8）：} 光谱噪声为平稳白噪声，相邻点差分稳健估计 $\hat{\sigma}_{\text{noise}}\approx0.012\%$。\\
\textbf{处置：} 分块 bootstrap（$N=200$）已处理可能的误差相关性。

\item \textbf{谱可分解为缓变背景加周期条纹（A9）：} 反射率谱 $R(\sigma)=\text{背景}(\sigma)+\text{条纹}(\sigma)$，背景以 7 阶多项式（poly7）近似，失配 $\sim0.1\%$ 量级计入残差。\\
\textbf{依据：} 数据实验（EXP-5）三方案对比，poly7 零伪峰、跨角度极差最小，频域高通方案因 Gibbs 振铃触发红线出局。
\end{enumerate}